\documentclass{IEEEtran}
\usepackage[utf8]{inputenc}
\usepackage[T1]{fontenc}
\usepackage{geometry}
\usepackage{graphicx}
\usepackage{float}
\usepackage{hyperref}
\usepackage[english,indonesian]{babel}

\title{Optimasi Pemrosesan Template LaTeX Menggunakan Pendekatan Berbasis AI}
\author{Academic Support Agent}
\date{\today}

\begin{document}
\maketitle

\begin{abstract}

\end{abstract}

\chapter{PENDAHULUAN}


\chapter{TINJAUAN PUSTAKA}


% Original content (commented out):
% \documentclass[9pt, compsoc, technote, a4paper]{IEEEtran}
\IEEEoverridecommandlockouts
% The preceding line is only needed to identify funding in the first footnote. If that is unneeded, please comment it out.
\usepackage{cite}
\usepackage{amsmath,amssymb,amsfonts}
\usepackage{algorithmic}
\usepackage{graphicx}
\usepackage{textcomp}
\usepackage{xcolor}
\usepackage{multirow}
\usepackage{lipsum}

\usepackage{pgfplots}
\pgfplotsset{width=0.47\textwidth,height=0.33\textwidth,compat=1.9}
% \usepgfplotslibrary{external}   % Externalizing the figures
% \tikzexternalize

\def\BibTeX{{\rm B\kern-.05em{\sc i\kern-.025em b}\kern-.08em
    T\kern-.1667em\lower.7ex\hbox{E}\kern-.125emX}}
    
\begin{document}

% Author and Titling
\title{Internship Title
\thanks{Acknowledge any supporting organization or department (excluding the internship company). If none, comment out or delete this line.}
}

\author{\small{Firstname Lastname$^{1}$} \\ \footnotesize{Supervised by Firstname Lastname$^{2}$}
\thanks{$^{1}$School of Science and Engineering, University of Tsukuba.}%
\thanks{$^{2}$Supervisor Position Title, Company Name.}
}

\maketitle

\begin{abstract}
This report provides an unofficial template for an internship report for Bachelor's Program in Interdisciplinary Engineering, University of Tsukuba. IDE guidelines call for a minimum length of five typed pages (A4). The IEEEtran LaTeX class with 9pt fontsize, compsoc, and technote options were used, with minor modifications. The sections included in this template are only suggestions; authors are free and encouraged to structure reports in a manner which best suits the needs of the topic and nature of the internship. As of June 2, 2021, the IDE steering committee's requirement is five typed A4 pages detailing daily activities at the internship, in addition 
\end{abstract}

\begin{IEEEkeywords}
(optional) List keywords that pertain to your internship.
\end{IEEEkeywords}

\section{Introduction}
\lipsum[1]

\section{Background}
\lipsum[1]

\subsection{Figures}
This is an example of a subsection and inserting figures. Please use linewidth (with or without modifier from 0.0 to 1.0) to regulate the size of the figure. Alternatively, the height can be regulated with height={value}. Beware that defining both width and height will sacrifice the figure's aspect ratio. For aesthetics, [b], [t], [h], or a combination of these options can be used to place figures at the bottom, top, or 'here', respectively.

\begin{figure}[t]
    \centering
    \includegraphics[width=0.5\linewidth]{fig/Picture1.png}
    \caption{\textit{Reprinted from [cite-here] (Remove if figure is original work)} Insert figure description here}
    \label{figure_1}
\end{figure}

\subsubsection{Subsubsection Example}
This is an example of a subsubsection.

\subsection{Citations and References}
Citations and references are not required, but in-text citations in IEEE-style is supported through the BibLaTeX package. Bibliography entries can be made in the refs.bib file. It can be cited in-text in this manner \cite{example_1}.

\section{Internship Activities}
\subsection{yyyy-mm-dd}
\begin{itemize}
    \item Example of daily internship activity for this day.
\end{itemize}

\subsection{yyyy-mm-dd}
\begin{itemize}
    \item Example of daily internship activity for this day.
\end{itemize}

\subsection{yyyy-mm-dd}
\begin{itemize}
    \item Example of daily internship activity for this day.
\end{itemize}

\section{Conclusion}
\lipsum[1]

\section*{Acknowledgment}
(Optional) Acknowledgements for any individual or organizations can be made here. Monetary support from grants or support programs can also be acknowledged in the footnote thanks section at the start of the document.

\bibliographystyle{ieeetr}   % We choose the "ieeetr" reference style
\bibliography{IEEEabrv,refs}          % Entries are in the refs.bib file

\end{document}


\end{document}